\begin{figure}[htbp]\centering%
\xeCJKDeclareCharClass{CJK}{"2460->"2473}%
\resizebox{\linewidth}{!}{%
\begin{tikzpicture}[x=1cm,y=1cm,line width=0.6pt,font=\small,
  grp/.style={draw=black!80,fill=black!5},
  box/.style={draw=black!80,fill=white,align=center,minimum height=0.9cm,inner sep=3pt},
  lab/.style={font=\footnotesize,align=center,inner sep=1.5pt},
  arr/.style={-{Stealth[length=2mm]},draw=black!80}]
% console
\node[box,fill=black!5,minimum width=14.9cm,minimum height=0.8cm] (ui) at (7.5,9.2) {操作员控制台（下发任务 · 人工审批 · 查看溯源图）};
% left group
\draw[grp] (0,1.7) rectangle (3.9,7.8);
\node[align=center,font=\small] at (1.95,7.35) {自研演示智能体\\（百炼大模型）};
\node[box,minimum width=2.1cm] (task) at (1.25,6.0) {任务输入};
\node[box,minimum width=2.1cm,text width=1.9cm] (llm) at (1.25,3.6) {大模型调用\\/ 子 Agent};
\draw[arr] (task)--(llm);
\node[box,dashed,minimum width=1.0cm,text width=0.9cm,minimum height=2.6cm,font=\footnotesize] (sdk) at (3.1,5.4) {溯界\\SDK\\（宿主\\Hook\\接口）};
\node[lab,text=black!70] at (1.95,1.15) {可替换为 Claude Code / Cursor\\等宿主（需接入 SDK）};
% gateway
\draw[grp] (5.6,0.4) rectangle (9.8,7.8);
\node[font=\small] at (7.9,7.45) {溯界安全网关};
\node[lab] at (7.9,7.0) {SDK + MCP 中间代理};
\node[box,minimum width=3.4cm] (s1) at (7.9,6.0) {① 多源注入检测\\（直接 + 间接）};
\node[box,minimum width=3.4cm] (s2) at (7.9,4.6) {② 意图对齐研判};
\node[box,minimum width=3.4cm] (s3) at (7.9,3.2) {③ 策略决策\\（放行 / 审批 / 阻断）};
\node[box,minimum width=3.4cm] (s4) at (7.9,1.8) {④ 全链路溯源};
\draw[arr] (s1)--(s2); \draw[arr] (s2)--(s3); \draw[arr] (s3)--(s4);
% tools
\draw[grp] (11.4,2.6) rectangle (15,7.8);
\node[font=\small] at (13.2,7.35) {MCP 服务器与工具};
\node[box,minimum width=2.9cm] (t1) at (13.2,6.2) {文件 / Git};
\node[box,minimum width=2.9cm] (t2) at (13.2,5.0) {网页 / HTTP};
\node[box,minimum width=2.9cm] (t3) at (13.2,3.8) {工单 / 邮件 / 记忆};
% external
\node[box,draw=red!60!black,text=red!60!black,minimum width=3.6cm,text width=3.4cm] (ext) at (13.2,1.0) {外部不可信内容\\{\footnotesize 恶意网页 · 投毒文档\\恶意 MCP server}};
\draw[arr,draw=red!60!black] (ext.north)--(13.2,2.6) node[lab,midway,right,text=red!60!black] {注入载荷};
% arrows
\draw[arr] (0.55,8.8)--(0.55,6.45) node[lab,pos=0.25,right] {下发任务};
\draw[{Stealth[length=2mm]}-{Stealth[length=2mm]},draw=black!80] (5.9,8.8)--(5.9,3.2)--(s3.west);
\node[lab,left] at (5.9,8.3) {ASK 审批};
\draw[arr,dashed] (sdk.east|-0,6.0)--(5.6,6.0) node[lab,midway,above] {进程内调用\\上下文 /\\输出登记};
\draw[arr] (llm.east)--(5.6,3.6) node[lab,midway,above] {MCP 工具调用};
\draw[arr] (9.8,6.1)--(11.4,6.1) node[lab,midway,above] {放行后转发};
\draw[arr] (11.4,4.6)--(9.8,4.6) node[lab,midway,below] {结果回流\\检测};
\end{tikzpicture}}
\caption{溯界系统架构。保留原图功能分区；当前原型的对齐说明用于解释已有决策，实际执行顺序见正文。}\label{fig:architecture}
\end{figure}
